\section{Coefficient functions through NNLO QCD}\label{sec:coefs}

The strong coupling and the quark mass require the regular
\qcd\ renormalization according to
\begin{equation}
  \begin{split}
    g_0 = \left( \frac{\mu \, e^{\EulerGamma / 2}}{\sqrt{4 \pi}}
    \right)^\ep Z_g\,g\,,\qquad m_0 = Z_m\,m\,,
    \label{eq:msbar}
  \end{split}
\end{equation}
where we write the renormalization constants $Z_g$ and $Z_m$ as
\begin{equation}
  \begin{split}
    Z_g &=
    1 - \gcoup{2}\frac{\beta_0}{2\ep} +
    \gcoup{4}\left(\frac{3\beta_0^2}{8\ep^2}
    -\frac{\beta_1}{4\ep}\right)   + \order{g^6}\,,
    \\
    Z_m &=
    1
    -\gcoup{2}\frac{\gamma_{m,0}}{2\ep}
    +\gcoup{4}\left[\frac{1}{\ep^2}\left(
      \frac{\gamma_{m,0}^2}{8} + \frac{\beta_0\gamma_{m,0}}{4}\right)
      -\frac{\gamma_{m,1}}{4\ep}\right] + \order{g^6}\,,
    \label{eq:rgconsts}
  \end{split}
\end{equation}
with
\begin{equation}
  \begin{split}
      \beta_0 &= \frac{11}{3} C_A - \frac{4}{3} \TF\,, \qquad
      \beta_1 = \frac{34}{3} C_A^2 - \left( 4 C_F + \frac{20}{3} C_A
      \right) \TF \,,\\
      \gamma_{m,0} &= 6 C_F\,,\qquad
      \gamma_{m,1} = \frac{97}{3} C_A C_F + 3 C_F^2
      - \frac{20}{3} C_F \TF \,.
    \label{eq:anombg}
  \end{split}
\end{equation}
$C_F$ and $C_A$ are the quadratic Casimir eigenvalues of the fundamental
and the adjoint representation of the gauge group,
respectively. Furthermore, $\TF=Tn_F$, with $n_F$ the number of quark
flavors, and $T$ the trace normalization in the fundamental
representation. For SU($N_c$), it is $C_F=(N_c^2-1)/(2N_c)$, $C_A=N_c$,
and $T=1/2$.

In addition, the flowed quark fields also require renormalization
according to
\begin{equation}
  \begin{split}
    \chi_{f,R} = \sqrt{Z_\chi}\,\chi_{f}\,,
  \end{split}
\end{equation}
leading to the factor $Z_\chi$ in the definition of the operators
$\tcalo_{3,4}$ in \eqn{eq:tops}.  $Z_\chi$ differs from the quark-field
renormalization of regular \qcd.
Through \nlo, the \msbar\ result has been evaluated in
\citere{Luscher:2013cpa}. To determine the renormalization constant at
\nnlo\ as required by our calculation, we may use the fact that the
coefficient $c_3(t)$ must be finite after renormalization. Writing the
\msbar\ expression as
\begin{equation}
  \begin{split}
    Z^{-1}_\chi &=
    1
      -\gcoup{2}\frac{\gamma_{\chi,0}}{2\ep}
      +\gcoup{4}\left[\frac{1}{\ep^2}\left(
        \frac{\gamma_{\chi,0}^2}{8} +
        \frac{\beta_0\gamma_{\chi,0}}{4}\right)
        -\frac{\gamma_{\chi,1}}{4\ep}\right] + \order{g^6}\,,
    \label{eq:zchi}
  \end{split}
\end{equation}
we find\footnote{Note that since $c_4=0$ at leading
  order, this coefficient only requires \nlo\ renormalization.}
\begin{equation}
  \begin{split}
  \gamma_{\chi,0} &= 6C_F\,,\\
  \gamma_{\chi, 1} &=
  C_A C_F \left( \frac{223}{3} - 16 \ln 2 \right) - C_F^2 \left( 3 + 16
  \ln 2 \right) - \frac{44}{3} C_F \TF\,.
  \end{split}
\end{equation}
This allows us to evaluate the coefficients of the energy-momentum
tensor in the \msbar\ scheme through \nnlo\ \qcd:\footnote{For
  convenience of the reader, we provide the expressions for
  $c_1,\ldots,c_4$ also in electronic form in an ancillary file with
  this paper.}
\begin{equation}
  \begin{split}
    c_1(t) =& \, \frac{1}{g^2} \Bigg\{ 1 + \gcoup{2}
    \left[ - \frac{7}{3}  C_A + \frac{3}{2}
      \TF -  \beta _0 \, L (\mu, t) \right] \\
	& \quad + \gcoup{4} \Bigg[ - \beta_1 \, L
      (\mu, t) + C_A^2 \left( - \frac{14482}{405} - \frac{16546}{135}
      \ln 2 + \frac{1187}{10} \ln 3 \right) \\
      & \qquad + C_A \TF \Bigg( \frac{59}{9}
      \text{Li}_2\left(\frac{1}{4}\right) + \frac{10873}{810} +
      \frac{73}{54} \pi^2 - \frac{2773}{135} \ln 2 + \frac{302}{45} \ln
      3 \Bigg) \\
	& \qquad + C_F \TF \Bigg( - \frac{256}{9}
      \text{Li}_2\left(\frac{1}{4}\right) +\frac{2587}{108} -
      \frac{7}{9} \pi^2 - \frac{106}{9} \ln 2 - \frac{161}{18} \ln 3
      \Bigg) \Bigg] \\
    & \quad + \order{g^6} \Bigg\} \, ,
    \label{eq:c1}
  \end{split}
\end{equation}
\begin{equation}
	\begin{split}
          c_2(t) =& \, \frac{1}{4 g^2} \Bigg\{ - 1 +  \gcoup{2}
          \left[ \frac{25}{6}  C_A - 3  \TF + \beta _0 \,
      L (\mu, t) \right] \\
	& \quad +  \gcoup{4} \Bigg[ \beta _1 \, L (\mu, t) +
      C_A^2 \left(\frac{56713}{1620} - \frac{1187}{10} \ln 3 +
      \frac{16546}{135} \ln 2 \right) \\
	& \qquad + C_A  \TF \Bigg( - \frac{59}{9}
      \text{Li}_2\left(\frac{1}{4}\right) - \frac{6071}{405} -
      \frac{73}{54} \pi^2 + \frac{2287}{135} \ln 2 - \frac{361}{90} \ln
      3 \Bigg) \\
	& \qquad + C_F \TF \Bigg( \frac{220}{9}
      \text{Li}_2\left(\frac{1}{4}\right) - \frac{1757}{54} +
      \frac{10}{9} \pi ^2 - \frac{164}{9} \ln 2 + \frac{247}{9} \ln 3
      \Bigg) \Bigg] \\
	& \quad  + \order{g^6} \Bigg\} \, ,
    \label{eq:c2}
  \end{split}
\end{equation}
\begin{equation}
	\begin{split}
	  c_3 (t) =& \, \frac{1}{4} \Bigg\{1 + \gcoup{2}
          \left( \frac{3}{2} C_F +
          \frac{\gamma_{\chi, 0}}{2} L (\mu, t) \right) \\
			& \quad + \gcoup{4}
          \Bigg[ \frac{\gamma_{\chi, 0}}{4} \left(\beta _0+
            \frac{\gamma_{\chi, 0}}{2} \right) \Big( L^2 (\mu, t) + L
            (\mu, t) \Big) + \frac{\gamma_{\chi, 1}}{2} L (\mu, t) \\
			& \qquad + C_F^2 \Bigg( - \frac{137}{9}
            \text{Li}_2\left(\frac{1}{4}\right) - \frac{559}{216} +
            \frac{103}{108} \pi^2 - \frac{1736}{27} \ln 2 +
            \frac{122}{3} \ln 3 - 4 \ln^2 2 \Bigg) \\
			&  \qquad + C_F  \TF \Bigg( -
            \frac{136}{9}  \text{Li}_2\left(\frac{1}{4}\right) -
            \frac{3377}{810} - \frac{7}{9} \pi^2 + \frac{1232}{135} \ln
            2 - \frac{136}{15} \ln 3 \Bigg) \\
			& \qquad + C_A C_F \Bigg( - \frac{365}{9}
            \text{Li}_2 \left(\frac{1}{4}\right) + \frac{261829}{3240} +
            \frac{77}{108} \pi^2 + \frac{5788}{45} \ln 2 \\
			& \qquad \quad - \frac{2102}{15} \ln 3 - 4 \ln^2
            2 \Bigg) \Bigg] \\
			& \quad + \order{g^6} \Bigg\} \, ,
    \label{eq:c3}
	\end{split}
\end{equation}
\begin{equation}
  \begin{split}
    c_4 (t) =& \, \frac{C_F}{2} \, \Bigg\{ \gcoup{2} +
    \gcoup{4} \Bigg[ \left( \beta _0
      + \frac{\gamma_{\chi, 0}}{2} \right) L (\mu, t) \\
      & \qquad + C_F \Bigg( - \frac{161}{18}
      \text{Li}_2\left(\frac{1}{4}\right) - \frac{41}{54} -
      \frac{55}{108} \pi^2 - \frac{1105}{27} \ln 2 +
      \frac{101}{6} \ln 3 \Bigg) \\
      & \qquad + \TF \Bigg( \frac{25}{9}
      \text{Li}_2\left(\frac{1}{4}\right) -
      \frac{20573}{1620} + \frac{5}{18} \pi^2 +
      \frac{6559}{135} \ln 2 - \frac{679}{30} \ln 3 \Bigg)
      \\
      & \qquad + C_A \Bigg( \frac{257}{36}
      \text{Li}_2\left(\frac{1}{4}\right) - \frac{137}{405}
      + \frac{11}{216} \pi^2 - \frac{419}{90} \ln 2 +
      \frac{1157}{60} \ln 3 \Bigg) \Bigg] \\
    & \quad + \order{g^6} \Bigg\} \, ,
    \label{eq:c4}
  \end{split}
\end{equation}
where we introduced the parameter\footnote{This parameter is motivated
  by the product of the typical factor $(8\pi t)^{\ep}$ occurring in
  flow-time integrals\,\cite{Luscher:2010iy}, and the usual definition
  of the renormalization scale in the \msbar\ scheme, see
  \eqn{eq:msbar}: $(8\pi t)^\ep (\mu^2e^{\EulerGamma}/(4\pi))^\ep = 1 +
  \ep\,L(\mu,t) + \order{\ep^2}$.}
\begin{equation}
	L (\mu, t) \equiv \ln \left( 2 \mu^2 t \right) + \EulerGamma\,,
        \label{eq:lmut}
\end{equation}
with the Euler-Mascheroni constant $\EulerGamma=0.57721\ldots$. Through
\nlo, these results are in full agreement with those of
\citere{Suzuki:2013gza,Makino:2014taa}. We have carried out the
calculation in the general $R_\xi$ gauge of regular \qcd; the fact that
the gauge-parameter dependence cancels in the final result serves as
another welcome check. The gauge parameter $\kappa$ of \eqn{eq:flow} has
been set to 1.

While the energy-momentum tensor $T_{\mu\nu}(x)$ is
renormalization-scheme independent, this is not necessarily the case for
the operators $\tcalo_{i,\mu\nu}(t,x)$ and the coefficient functions
$c_i(t)$. Since $\tcalo_{1,\mu\nu}$ and $\tcalo_{2,\mu\nu}$ do not
require operator renormalization, their matrix elements as well as the
coefficient function are indeed renormalization-scheme independent.  On
the other hand, using the quark-field renormalization $Z_\chi$ of
\eqn{eq:zchi} in the \msbar\ scheme, matrix elements of
$\tcalo_{i,\mu\nu}$ and coefficient functions $c_i(t)$ become
explicitly dependent on the renormalization scale $\mu$ for
$i\in\{3,4\}$.

However, this renormalization-scheme dependence can be avoided by
introducing so-called ``ringed'' quark fields as suggested
in~\citere{Makino:2014taa}. This corresponds to replacing $Z_\chi$ in
\eqn{eq:tops} by
\begin{equation}
  \begin{split}
    \mathring{Z}_\chi(t) = \frac{-2 N_c \, n_F}{(4 \pi t)^2
      \langle\bar{\chi}_f(t,x)\overleftrightarrow
      {\slashed D}\chi_f(t,x)\rangle}\,.
    \label{eq:zf}
  \end{split}
\end{equation}
Currently, $\mathring{Z}_\chi(t)$ is available only through \nlo\ \qcd.
Its explicitly $\mu$-dependent terms can be reconstructed from the
requirement that $c_i(t)$ must be finite and $\mu$-independent for $i\in
\{3,4\}$ though. In this way we find for the ratio to the
\msbar\ quark-field renormalization constant of \eqn{eq:zchi}:
\begin{equation}
	\begin{split}
	  \zeta_\chi&\equiv \frac{\mathring{Z}_\chi}{Z_\chi} = 1 + \gcoup{2}
          \left( \frac{\gamma_{\chi, 0}}{2} L(\mu ,t) - 3 C_F \ln 3 - 4
          C_F \ln 2 \right) \\ &
                \qquad + \gcoup{4} \Bigg\{
                  \frac{\gamma_{\chi, 0}}{4} \left( \beta_0 +
                  \frac{\gamma_{\chi, 0}}{2} \right) L^2 (\mu, t) +
                  \Big[ \frac{\gamma_{\chi, 1}}{2} -
                  \frac{\gamma_{\chi, 0}}{2} \left( \beta_0 +
                  \frac{\gamma_{\chi, 0}}{2} \right) \ln 3 \\
                & \qquad \qquad -
                  \frac{2}{3}
                  \gamma_{\chi, 0} \left(
                  \beta_0 + \frac{\gamma_{\chi, 0}}{2}\right) \ln 2
                  \Big] L(\mu,t)
                  + C_2 \Bigg\} \\ &\qquad + \order{g^6} \,.
                \label{eq:zfchi}
	\end{split}
\end{equation}
The constant $C_2$ cannot be determined in this way, but requires a
dedicated \three-loop calculation of the \two-point function occurring
in the denominator of \eqn{eq:zf}. A detailed outline of this
calculation is beyond the scope of this paper; it will be presented
together with a more complete description of our setup in a forthcoming
publication\,\cite{Artz:2019bpr}. At this point, we simply quote the
numerical value of this result up to three significant digits, which is
more than sufficient in the light of the theoretical uncertainties to be
discussed below:\footnote{Note that the calculation of
  $\mathring{Z}_\chi$ also provided an independent check for $Z_\chi$.}
\begin{equation}
  C_2 = -23.8 \, C_A C_F + 30.4 \, C_F^2 - 3.92 \, C_F T_F\,.
\end{equation}
Multiplication of $c_3(t)$ and
$c_4(t)$ in Eqs.\,(\ref{eq:c3}) and (\ref{eq:c4}) by this ratio makes
also these coefficients formally $\mu$-independent, i.e.,
\begin{equation}
  \begin{split}
    \mu\dderiv{}{\mu}\{c_1,c_2,\ringc_3,\ringc_4\}=0\,,\quad
    \text{where}\quad \ringc_i\equiv \zeta_\chi^{-1} c_i\,.
  \end{split}
\end{equation}
As in any perturbative calculation, the $\mu$-independence only holds up
to higher orders in $g$. The decrease of the residual $\mu$-dependence
is thus commonly used as a qualitative check of the perturbation
expansion for the specific observable under consideration. We thus study
the $\mu$-dependence of the four coefficients after dividing $c_3$
and $c_4$ by the ratio $\zeta_\chi$ defined in \eqn{eq:zfchi}. We fix a
characteristic value for the flow time $t$ and vary the renormalization
scale $\mu$ around the central value $\mu_0$, which we define such that
$L(\mu_0,t)=0$, cf.\ \eqn{eq:lmut}, i.e.
\begin{equation}
  \begin{split}
    \mu_0 = \frac{e^{-\EulerGamma/2}}{\sqrt{2t}}\,.
    \label{eq:mu0}
  \end{split}
\end{equation}
Figures\,\ref{fig:mu3gev} and \ref{fig:mu130gev} show the leading order
(\lo), \nlo, and the \nnlo\ approximation of $c_1$, $c_2$, $\ringc_3$,
and $\ringc_4$ as functions of the renormalization scale for two
different values of the flow time $t$, corresponding to $\mu_0=3$\,GeV
and $\mu_0=130$\,GeV, respectively. In the former case, we set $n_F=3$,
in the latter $n_F=5$. We use $\alpha^{(n_F=5)}_s(M_Z)=0.118$ in order
to evaluate the input values for the couplings,
$g^{(n_F=3)}(3\,\text{GeV})=1.77$ and
$g^{(n_F=5)}(130\,\text{GeV})=1.19$.  The $\mu$-variation of the strong
coupling constant $g(\mu)$ is determined by numerically solving the
corresponding renormalization-group equation with the help of
\texttt{RunDec}\,\cite{Chetyrkin:2000yt,Herren:2017osy} at \one-, \two-,
and \three-loop level for the \lo, \nlo, and the \nnlo\ curve,
respectively.  In \fig{fig:mu3gev}, the value of $t$ is chosen such that
the central scale of \eqn{eq:mu0} is $\mu_0=3$\,GeV. At this central
scale, the \nnlo\ corrections increase the modulus of the coefficients $c_1$ and $c_2$ by 10\% and 13\%
relative to NLO, respectively. This is within twice the
\nlo\ uncertainty due to missing higher-order effects as estimated by
varying $\mu/\mu_0$ between 1/2 and 2, where one finds $7.3$\% for
$c_1$, and $8.0\%$ for $c_2$. We are therefore confident that the
\nnlo\ uncertainty estimated in the same way is rather reliable: it is
given by $5.7$\% for $c_1$ and $7.2$\% for $c_2$. Note that the dominant
contribution to these numbers comes from the downward variation of
$\mu$, where $g(\mu)$ starts to become sensitive to the non-perturbative
region. The behavior of $c_1$ and $c_2$ towards larger values of $\mu$
seems to suggest that this uncertainty estimate may actually be too
conservative.


\begin{figure}
  \begin{center}
    \begin{tabular}{cc}
      \includegraphics[width=.45\textwidth,bb=20 0 300 220]{images/c1_3.pdf} &
      \includegraphics[width=.45\textwidth,bb=20 0 300 220]{images/c2_3.pdf} \\
      \includegraphics[width=.45\textwidth,bb=20 0 300 220]{images/c3_3.pdf} &
      \includegraphics[width=.45\textwidth,bb=20 0 300 220]{images/c4_3.pdf}
    \end{tabular}
  \end{center}
	\caption{Renormalization-scale dependence of the coefficients
          $c_1$, $c_2$, $\ringc_3=\zeta^{-1}_\chi c_3$,
          $\ringc_4=\zeta^{-1}_\chi c_4$, defined in
          \eqs{eq:c1}--\noeqn{eq:c4}, with $\zeta_\chi$ from
          \eqn{eq:zfchi}. The dotted black, dashed blue, and solid red
          curve correspond to keeping terms up to order $(g^2)^{n-1}$ in
          $c_1$ and $c_2$, and $(g^2)^{n}$ in $\ringc_3$ and
          $\ringc_4$, with $n=0,1,2$, respectively. The central
          scale is set to $\mu_0=3$\,GeV, corresponding to $t = 3.12
          \cdot 10^{-2} /\text{GeV}^2$, see \eqn{eq:mu0}.
          The number of flavors is set to $n_F=3$.}
	\label{fig:mu3gev}
\end{figure}


\begin{figure}
\begin{center}
    \begin{tabular}{cc}
      \includegraphics[width=.45\textwidth,bb=20 0 300 220]{images/c1_130.pdf} &
      \includegraphics[width=.45\textwidth,bb=20 0 300 220]{images/c2_130.pdf} \\
      \includegraphics[width=.45\textwidth,bb=20 0 300 220]{images/c3_130.pdf} &
      \includegraphics[width=.45\textwidth,bb=20 0 300 220]{images/c4_130.pdf}
    \end{tabular}
  \end{center}
	\caption{Same as \fig{fig:mu3gev}, but for $\mu_0=130$\,GeV (or
          $t=1.66\cdot 10^{-5} /\text{GeV}^2$), and $n_F=5$.}
	\label{fig:mu130gev}
\end{figure}


As opposed to the gluonic coefficients $c_1$ and $c_2$, the coefficients
of the fermionic operators $c_3$ and $c_4$ exhibit a residual scale
dependence only from the \nlo\ term onwards. One therefore expects a
stronger $\mu$-dependence at \nnlo\ for these terms. Nevertheless, for
$\ringc_3$, the estimate of the theory uncertainty due to scale
variation still decreases from $9.8$\% to $8.1$\%. The increase of the
result due to the \nnlo\ effects is $8.3$\% relative to the \nlo\ result
at $\mu=\mu_0$.

The behavior of $\ringc_4$, on the other hand, is less
satisfactory at $\mu_0=3$\,GeV. The \nnlo\ effects more than double the
\nlo\ result in this case, and the uncertainty estimate due to scale
variation actually \textit{increases} from 47\% to 71\% when going from
\nlo\ to \nnlo.  Note, however, that $c_4=0$ at \lo, which means that
this coefficient is numerically sub-dominant.

As one would expect, for $\mu_0=130$\,GeV, the perturbative behavior of
all coefficients is significantly improved, cf.\ \fig{fig:mu130gev}. For
$c_1$, $c_2$, and $\ringc_3$, the scale uncertainty is at the
sub-percent level already at \nlo; at \nnlo, it amounts to less than
$0.2$\% in all three cases. The effect of the \nnlo\ corrections
relative to the \nlo\ result is about $2$\% for $c_1$ and $c_2$, and
$0.8$\% for $\ringc_3$. Also for $\ringc_4$, the
situation improves significantly: the \nnlo\ terms add 38\% to the
\nlo\ result, and the uncertainty goes down from 10\% at \nlo\ to
$5.8$\% at \nnlo.

It is also worth pointing out that the choice of the central scale
$\mu_0$ as defined in \eqn{eq:mu0} seems justified by the behavior of
the successively higher orders. In almost all cases, the \nlo\ and the
\nnlo\ corrections are \textit{both} relatively small at $\mu=\mu_0$. At
the same time, the \nnlo\ corrections relative to the \nlo\ result are
always smaller than the \nlo\ corrections compared to the
\lo\ result. The only exception to this is $\ringc_4$ at $\mu_0=3$\,GeV,
where, however, no choice of $\mu$ seems to stand out over any other.


In summary, we conclude that the \nnlo\ terms lead to a significant
improvement of the perturbative accuracy of the Wilson coefficients.
